% Szkielet dla pracy inżynierskiej pisanej w języku angielskim.

\documentclass[english,bachelor,a4paper,oneside,draft]{ppfcmthesis}

\usepackage[utf8]{inputenc}
\usepackage[OT4]{fontenc}

% Authors here.
\author{%
   Jan Kończak \album{84822} \and 
   Tomasz Żurkowski \album{84915}}
\authortitle{}                                        % Do not change.
%\title{\Huge JPaxos \normalsize State~machine~replication~based~on~the~Paxos~protocol}
 \title{JPaxos State~machine~replication based~on~the~Paxos~protocol}
\ppsupervisor{dr hab. inż. Paweł T. Wojciechowski} % Your supervisor comes here.
\ppyear{2011}                                         % Year of final submission (not graduation!)

\begin{document}

% Front matter starts here
\frontmatter\pagestyle{empty}%
\maketitle\cleardoublepage%

% Blank info page for "karta dyplomowa"
\thispagestyle{empty}\vspace*{\fill}%
\begin{center}Tutaj przychodzi karta pracy dyplomowej;\\oryginał wstawiamy do wersji dla archiwum PP, w pozostałych kopiach wstawiamy ksero.\end{center}%
\vfill\cleardoublepage%

% Table of contents.
\pagenumbering{Roman}\pagestyle{ppfcmthesis}%
\tableofcontents* \cleardoublepage%

% Main content of your thesis starts here.
\mainmatter%

\chapter{Introduction}

With the increase of usage and importance of Internet services in everyday life increases also demand on reliability and constant availability.
Providing these features to the services is a complex matter, often requiring various trade-offs between user satisfaction and service resource demands.

Most common and effective method of providing reliable services is replication. Applying the replication requires only minor changes for the non-replicated service to become a replicated ones, what makes it usable both for creating new services following the well known good patterns as well as transforming existing services into replicated one.

Each replication approach relies on certain group communication mechanism -- usually providing strong guarantees. Total-order (atomic) broadcast is a very useful for interconnecting replicas -- it simplifies the design and implementation of replicated protocols, as it allows for dividing the system into two main modules with clearly visible responsibilities.

The services vary, and so the approaches for replication must also vary. A library able to replicate every service would not only be hard to implement, but also could be too complicated for the user.

We chose to implement one of the replication approaches -- an approach we trust to have bright future.
We implement active replication of services -- each copy does the same work. The replicas must work identically, therefore we require the service to behave deterministically.

Every library for replicating services should address its main purpose -- guarantee high availability of service as well as durability of the service work results -- while having minimum requirements for the environment.

Our software design goal is to work well in real networks, survive temporary network crashes and be able to recover from replica crashes as well.

{
\bfseries
The goal of this bachelor's thesis is to create Java library implementing state machine replication primitive. The library should support the crash-recovery model of failures and should tolerate message loss and communication delays.
}

\section{Thesis organisation}
Second part of this chapter characterizes shortly articles concerning Paxos and present most significant definitions as well as theoretical results concerning Paxos.

The second chapter provides description of Paxos algorithm used by us, including the modifications and implementation details we used.

In chapter 3 we describe the design of JPaxos -- its architecture, modules, important data structures and threading approach.

Implemented features, chapter 4, is devoted to the implementation work related to the state machine replication, including the snapshotting, catch-up and proper achieving transparency for the service.

Fifth chapter describes the recovery related issues, both for Paxos and for state machine replication.

Last chapter concludes the thesis. %TODO: finish as this will be written

\section{Related work}

The Paxos and MultiPaxos algorithm have been first described by L. Lamport in \textit{The Part-Time Parliament}~\cite{Lam98}. Since the article presented the algorithm in quite specific form, the Paxos protocol has been described again by Lamport in \textit{Paxos Made Simple}~\cite{Lam01}.

Since then, for over 7 years, the Paxos protocol has been described from theoretical point of view. Improvements and their proofs of correctness were presented as well as behaviour with byzantine crashes has been characterised.

First Paxos protocol uses have been described in 2006 and 2007 (\textit{Paxos made\linebreak live}~\cite{CGR07}). The articles describe certain implementation issues concerning the Paxos algorithm in Google Chubby distributed filesystem.

Year later, in \textit{Paxos for system builders}~\cite{AK08}, the Paxos protocol has been described from programmer point of view.

In 2009, and subsequently in 2010, M. Primi et alii released two implementations of Paxos designed to achieve good performance. Both of these are released as open-source. The first -- libPaxos$^2$ -- is described in his master thesis, the latter -- RingPaxos -- is described in \textit{Ring Paxos: A High-Throughput Atomic Broadcast Protocol}~\cite{Mar10}.

The crash-recovery has been also the subject among the articles. The \textit{Atomic Broadcast in Asynchronous Crash-Recovery Distributed Systems}~\cite{rodriguez2000atomic} provides definitions concerning recovery as well as shows method to transform any crash-stop consensus protocol to crash-recovery one.

A summary of the crash-recovery approaches for Paxos algorithm, \textit{Recovery on the Paxos Protocol}~\cite{Nun10}, has been written by N. Santos recently. All algorithms implemented by us are described there.

\section{Quid quis scripsit}
%FIXME:  ^^^^ ^^^^ ^^^^^^^^

Jan Kończak for the purpose of this thesis has prepared design and implementation of:
\begin{tightList}
  \item[\textbullet] catch-up module (Section \ref{sec:catch_up}),
  \item[\textbullet] snapshotting module (Section \ref{sec:snapshotting}),
  \item[\textbullet] service proxy  (Section \ref{sec:serviceProxy}),
  \item[\textbullet] benchmark project.
\end{tightList}

\noindent Tomasz Żurkowski has prepared design and implementation of:
\begin{tightList}
  \item[\textbullet] unit tests for JPaxos project,
  \item[\textbullet] client manager module, 
  \item[\textbullet] network module for communication between replicas,
  \item[\textbullet] view based recovery (Section \ref{sec:view_ss}),
  \item[\textbullet] disc writers for full stable storage recovery.
\end{tightList}

\noindent The remaining modules were result of our common work, and later division of work is just infeasible.

\chapter{Theory}

%Rozdział teoretyczny --- przegląd literatury naświetlający stan wiedzy na dany temat. 
%
%Przegląd literatury naświetlający stan wiedzy na dany temat obejmuje rozdziały pisane na podstawie
%literatury, której wykaz zamieszczany jest w części pracy pt.~\emph{Literatura} (lub inaczej \emph{Bibliografia},
%\emph{Piśmiennictwo}). W tekście pracy muszą wystąpić odwołania do wszystkich pozycji zamieszczonych w
%wykazie literatury. \textbf{Nie należy odnośników do literatury umieszczać w stopce strony.} Student jest
%bezwzględnie zobowiązany do wskazywania źródeł pochodzenia informacji przedstawianych w pracy,
%dotyczy to również rysunków, tabel, fragmentów kodu źródłowego programów itd. Należy także podać
%adresy stron internetowych w przypadku źródeł pochodzących z Internetu.

In this chapter the background for paxos algorithm is presented. First the consensus problem is described, th


Co to consensus, snapshot, service, catchup, deterministic state machinme, state machine, 
crash models, atomic broadcast, nierozwiazywalnosc paxos'a

\section{Consensus}

Consensus is a problem in distributed computing that encapsulates the task of group agreement in the presence of faluts.

\begin{description}
    \item[Validity] any value decided is a value proposed by some process,
    \item[Agreement] no two processes decide differently,
    \item[Termination] every correct process eventually decides,
    \item[Integrity] no process decides twice
\end{description}



\chapter{JPaxos}

\section{System model}

We assume that we replicate N servers - called processes in later discussion.
Servers communicate with each other using asynchronous communication with no
guarantees - messages can be lost and duplicated. Each server is running the
same state machine which is executing all requests from the clients.

\section{Client}

\section{Architecture}

\section{Paxos}
Proposer, acceptor, learner.

\section{Batching}

\section{Concurent instanses}

\section{Catch-up mechanism}

\section{Snapshoting}

\section{Storage}

\section{Client manager}

\section{Recovery}

\subsection{Crash Stop}

\subsection{Crash-recovery with stable storage}

\subsection{View based recovery}

\subsection{Epoch based recovery}

\section{Application}
MapService, DigestService - validation

PaxosSTM - []

%Rozdziały dokumentujące pracę własną studenta: opisujące ideę, sposób lub metodę 
%rozwiązania postawionego problemu oraz rozdziały opisujące techniczną stronę rozwiązania 
%--- dokumentacja techniczna, przeprowadzone testy, badania i uzyskane wyniki. 
%
%Praca musi zawierać elementy pracy własnej autora adekwatne do jego wiedzy praktycznej uzyskanej w
%okresie studiów. Za pracę własną autora można uznać np.: stworzenie aplikacji informatycznej lub jej
%fragmentu, zaproponowanie algorytmu rozwiązania problemu szczegółowego, przedstawienie projektu 
%np.~systemu informatycznego lub sieci komputerowej, analizę i ocenę nowych technologii lub rozwiązań
%informatycznych wykorzystywanych w przedsiębiorstwach, itp. 
%
%Autor powinien zadbać o właściwą dokumentację pracy własnej obejmującą specyfikację założeń i 
%sposób realizacji poszczególnych zadań
%wraz z ich oceną i opisem napotkanych problemów. W przypadku prac o charakterze 
%projektowo-implementacyjnym, ta część pracy jest zastępowana dokumentacją techniczną i użytkową systemu. 
%
%W pracy \textbf{nie należy zamieszczać całego kodu źródłowego} opracowanych programów. Kod źródłowy napisanych
%programów, wszelkie oprogramowanie wytworzone i wykorzystane w pracy, wyniki przeprowadzonych
%eksperymentów powinny być umieszczone na płycie CD, stanowiącej dodatek do pracy.
%
%\section*{Styl tekstu}
%
%Należy\footnote{Uwagi o stylu pochodzą częściowo ze stron Macieja Drozdowskiego~\cite{mdro}.} 
%stosować formę bezosobową, tj.~\emph{w pracy rozważono ......, 
%w ramach pracy zaprojektowano ....}, a nie: \emph{w pracy rozważyłem, w ramach pracy zaprojektowałem}. 
%Odwołania do wcześniejszych fragmentów tekstu powinny mieć następującą postać: ,,Jak wspomniano wcześniej, ....'', 
%,,Jak wykazano powyżej ....''. Należy unikać długich zdań. 
%
%,,Ilość'' i ,,liczba''. Proszę zauważyć, liczba dotyczy rzeczy policzalnych, np.~liczba osób, liczba zadań, procesorów. 
%Ilość dotyczy rzeczy niepoliczalnych, np.~ilość wody, energii. Należy starać się wyrażać precyzyjnie, tj.~zgodnie 
%z naturą liczonych obiektów.\footnote{(DW) Według wytycznych Rady Języka Polskiego obie formy są dopuszczalne
%zarówno do obiektów policzalnych, jak i niepoliczalnych. W tekstach technicznych warto być jednak precyzyjnym.}
%
%Niedopuszczalne są zwroty używane w języku potocznym. W pracy należy używać terminologii informatycznej, która ma 
%sprecyzowaną treść i znaczenie. Nie należy używać ,,gazetowych'' określeń typu: 
%silnik bazy danych, silnik programu, maszyna skryptowa, elektroniczny mechanizm, mapowanie, string, gdyż nie wiadomo 
%co one właściwie oznaczają. 
%
%Niedopuszczalne jest pisanie pracy metodą \emph{cut\&paste}, bo jest to plagiat i dowód intelektualnej indolencji autora.
%Dane zagadnienie należy opisać własnymi słowami. Zawsze trzeba powołać się na zewnętrzne źródła. 


\chapter{Conclusions}

In this thesis, we have studied the problem of state machine replication based
on Paxos algorithm defined in \cite{Lam98}. We described challenges faced in
implementation phase to make the library ready to use in production systems
with support of crash-recovery model, message loss and communication delays.

We first presented the Paxos algorithm with theoretical background. Then we
focused on problems which occurs when transforming one page pseudocode of
distributed algorithm to few thousands line of production code. We presented
the architecture of our library with short description of every module and data
structures used. The library has to handle unreliable network with message
loss and delays, so we presented the catch-up mechanism, which guarantees that
eventually every correct replica will be up to date. By introducing snapshot
mechanism, we solve a problem with infinite grow of log. To make the library
very efficient, we implemented several performance improvements like concurrent
instances, skipping redundant message and batching. To tolerate crash-recovery
model, four different recovery algorithms have been implemented and described in
this paper.

To test the library we implemented replicated hash map as an example of
deterministic service. We used it to verify that our library is working
correctly and that API we provide is user friendly. Our library is used as 
communication layer in PaxosSTM master thesis \cite{Tad10}.

Our future work includes further performance improvements to Paxos protocol and
implementing more crash-recovery algorithms with no use of stable storage.  It
would be also interesting to evaluate all recovery algorithms in network with
message loss and environment when failure of replica is a common scenario.
Finally, we would like to publish the library as open source project and test it
in production environment.

% TODO JK Acknowledgments (Wojciechowski, Nuno?, Andre?, Tadek i Maciej?)


% All appendices and extra material, if you have any.
\cleardoublepage\appendix%

\chapter{Appendix 1}

\section{Section 1}
\section{Section 2}


\chapter{Składanie dokumentu w systemie \LaTeX}

Po pierwsze to gratulacje --- dobry wybór. W tym rozdziale znajduje się
garść informacji o tym, jak poprawnie składać tekst pracy w systemie \LaTeX{} wraz z 
przykładami, które mają służyć do przeklejania do własnych dokumentów.

\section{Narzędzia}
Pracując pod systemem Windows, polecam:
\begin{itemize}
    \item MikTeX, \url{http://www.miktex.org/},
    \item JEdit, \url{http://www.jedit.org/},
    \item TeXlipse, \url{http://texlipse.sourceforge.net/},
    \item Kile, \url{http://kile.sourceforge.net/},
    \item Ghostview, Ghostscript (podgląd dokumentów PDF bez blokowania pliku):\\
        \url{http://www.cs.wisc.edu/~ghost/}. 
\end{itemize}

Po zainstalowaniu tych narzędzi wystarczy wykonać polecenie \texttt{compile.bat} (który
jest skryptem wsadowym dla Windows). Dla tych, którzy wolą nieco automatyzacji --- skrypt
\texttt{latexmk}, który jest w MikTeXu (a który potrzebuje zainstalowanego Perla) jest
również bardzo wygodny: \texttt{latexmk -pdf -pvc main.tex}.

\section{Edycja tekstu}
\chaptermark{Tytuł rozdziału, jeśli pełen się nie mieści\ldots{}}{}

\subsection{Struktura dokumentu}

Praca składa się z rozdziałów (\texttt{chapter}) i podrozdziałów (\texttt{section}).
Ewentualnie można również rozdziały zagnieżdzać (\texttt{subsection}, \texttt{subsubsection}),
jednak nie powinno się wykraczać poza drugi poziom hierarchii (czyli \texttt{subsubsection}).

\subsection{Akapity i znaki specjalne}

Każdy akapit to po prostu blok tekstu. Nieważne jak sformatowany --- to zrobi już
system $\LaTeX$.

Akapity rozdziela się od siebie przynajmniej jedną pustą linią. Podstawowe
instrukcje, które się przydają to \emph{wyróżnienie pewnych słów}. Można również
stosować \textbf{styl pogrubiony}, choć nie jest to generalnie zalecane.

Należy pamiętać o zasadach polskiej interpunkcji i ortografii. Po spójnikach 
jednoliterowych warto wstawić znak tyldy ($\sim$), który jest tak zwaną
,,twardą spacją'' i powoduje, że wyrazy nią połączone nie będą rozdzielane
na dwie linie tekstu.

Polskie znaki interpunkcyjne różnią się nieco od angielskich: to jest ,,polski'', a to jest
``angielski''. W kodzie źródłowym tego tekstu będzie widać różnicę.

Proszę również zwrócić uwagę na znak myślnika, który może być pauzą ,,---'' lub
półpauzą: ,,--''. Należy stosować je konsekwentnie. Do łączenia wyrazów używamy
zwykłego ,,-'' (\emph{północno-wschodni}), do myślników --- pauzy lub półpauzy.
Inne zasady interpunkcji i typografii można znaleźć w słownikach.

\subsection{Wypunktowania}

Wypunktowanie z cyframi:
\begin{enumerate}
    \item to jest punkt,
    \item i to jest punkt,
    \item a to jest ostatni punkt.
\end{enumerate}

\noindent
Po wypunktowaniach czasem nie warto wstawiać wcięcia akapitowego. Wtedy przydatne jest
polecenie \texttt{noindent}. Wypunktowanie z kropkami (tzw.~\emph{bullet list}) wygląda tak:
\begin{itemize}
    \item to jest punkt,
    \item i to jest punkt,
    \item a to jest ostatni punkt.
\end{itemize}

\noindent
Wypunktowania opisowe właściwie niewiele się różnią:
\begin{description}
    \item[elementA] to jest opis,
    \item[elementB] i to jest opis,
    \item[elementC] a to jest ostatni opis.
\end{description}


\subsection{Polecenia pakietu \texttt{ppfcmthesis}}

Parę poleceń zostało zdefiniowanych aby uspójnić styl pracy. Są one przedstawione poniżej
(oczywiście nie trzeba się do nich stosować).

\paragraph{Makra zdefiniowane dla języka angielskiego.} Są nimi: \texttt{termdef} oraz \texttt{acronym}.
Przykłady poniżej obrazują ich przewidywane użycie w tekście.
\begin{center}\footnotesize%
\begin{tabular}{l >{\rightskip\fill}p{12cm}}
\toprule
źródło   & \texttt{we call this a $\backslash$termdef\{Database Management System\} ($\backslash$acronym\{DBMS\})} \\ \cmidrule(lr){2-2}
docelowo & we call this a \termdef{Database Management System} (\acronym{DBMS}) \\ 
\bottomrule
\end{tabular}
\end{center}

\paragraph{Makra zdefiniowane dla języka polskiego.} Podobnie jak dla języka angielskiego zdefiniowano
odpowiedniki polskie: \texttt{defini\-cja}, \texttt{akronim} oraz \texttt{english} dla tłumaczeń angielskich
terminów. Przykłady poniżej obrazują ich przewidywane użycie w tekście.
\begin{center}\footnotesize%
\begin{tabular}{l >{\rightskip\fill}p{12cm}}
\toprule
źródło   & \texttt{nazywamy go $\backslash$definicja\{systemem zarządzania bazą danych\} ($\backslash$akronim\{DBMS\}, $\backslash$english\{Database Management System\})} \\ \cmidrule(lr){2-2}
docelowo & nazywamy go \definicja{systemem zarządzania bazą danych} (\akronim{DBMS}, \english{Database Management System}) \\ \bottomrule
\end{tabular}
\end{center}


\subsection{Rysunki}

Format wstawianych rysunków zależy od tego czy używa się do kompilacji polecenia
\texttt{latex}, czy też \texttt{pdflatex}. Oba powinny dać dokładnie ten sam wynik końcowy,
ale praca z nimi jest nieco inna.

\begin{description}
    \item[latex] To polecenie kompiluje źródła \LaTeX{}owe do pliku 
        z rozszerzeniem \texttt{dvi}. Ten plik można przeglądać przy pomocy specjalizowanych programów
        takich jak przykładowo Yap obecny z dystrybucją Mik\TeX{}a. Aby uzyskać docelowy plik \akronim{PDF}
        należy przekonwertować plik \texttt{dvi} przy pomocy programu \texttt{dvipdfm}. 
        
        \textbf{UWAGA:} korzystając z programu \texttt{latex}, wszystkie rysunki muszą być w formacie \akronim{EPS}
        (\english{encapsulated postscript}).

    \item[pdflatex] To polecenie kompiluje źródła \LaTeX{}owe bezpośrednio do pliku \akronim{PDF}.
    
        \textbf{UWAGA:} korzystając z programu \texttt{pdflatex}, wszystkie rysunki muszą być w formacie \akronim{PDF},
        \akronim{JPG} lub \akronim{PNG}.
\end{description}

Można oczywiście używać obu systemów --- wtedy pliki rysunków muszą po prostu być dostępne w obu formatach.

Wszystkie rysunki (w tym również diagramy, szkice i inne) osadzamy w środowisku 
\texttt{figure} i umieszczamy podpis \emph{pod} rysunkiem, w formie elementu \texttt{caption}. Rysunki powinny
zostać umieszczone u góry strony (osadzone bezpośrednio w treści strony zwykle utrudniają czytanie tekstu).
Rysunek~\ref{rys:plama} zawiera przykład pełnego osadzenia rysunku na stronie.

\begin{figure}[t] % możliwe opcje to 't' - top, 'b' - bottom, 'h' - 'here', ale zaleca się 't'
\centering\includegraphics[width=5cm]{figures/template/logo-pp}
\caption{Logo Politechniki Poznańskiej.}\label{rys:plama}
\end{figure}

\begin{figure}[t]
\centering\includegraphics[width=5cm]{figures/template/logo-pp}
\fcmfcaption{Logo Politechniki Poznańskiej. Formatowanie zgodne z wytycznymi FCMu.}\label{rys:plama2}
\end{figure}

Zasady FCMu sugerują nieco inne nagłówki rysunków. Dostepne są one poleceniem \texttt{fcmfcaption} (zob.~rysunek
\ref{rys:plama2}), jeśli ktoś woli mieć podpisy niespójne z rysunkami\ldots

\subsection{Tablice}

Tablice to piękna rzecz, choć akurat ich umiejętne tworzenie w \LaTeX{}u nie jest łatwe. 
Jeśli tablica jest skomplikowana, to pewnie łatwiej będzie ją wykonać w programie
OpenOffice, a następnie wyeksportować jako plik \akronim{PDF}. W każdym przypadku tablice wstawia się podobnie
jak rysunki, tylko że w środowisko \texttt{table}. Tradycja typograficzna sugeruje umieszczenie opisu tablicy, a więc
elementu \texttt{caption} ponad jej treścią (inaczej niż przy rysunkach).  

Tablica~\ref{tab:tabela} pokazuje pełen przykład.

\begin{table}[h]
\caption{Przykładowa tabela. Styl opisu jest zgodny z rysunkami.}\label{tab:tabela}
\centering\footnotesize%
\begin{tabular}{l c}
\toprule
artykuł & cena [zł] \\
\midrule
bułka   & $0,4$ \\
masło   & $2,5$ \\
\bottomrule
\end{tabular}
\end{table}

Zasady FCMu sugerują nieco inne nagłówki tablic. Dostepne są one poleceniem \texttt{fcmtcaption} (zob.~tablicę
\ref{tab:tabela2}), jeśli ktoś woli mieć podpisy niespójne z rysunkami\ldots

\begin{table}[h]
\fcmtcaption{Przykładowa tabela. Styl opisu jest zgodny z wytycznymi FCMu.}\label{tab:tabela2}
\centering\footnotesize%
\begin{tabular}{l c}
\toprule
artykuł & cena [zł] \\
\midrule
bułka   & $0,4$ \\
masło   & $2,5$ \\
\bottomrule
\end{tabular}
\end{table}


\subsection{Checklista}

W katalogu źródeł stylu \texttt{ppfcmthesis} znajduje się plik \texttt{CHECKLIST} --- należy
sprawdzić, czy nie popełniło się któregoś z wymienionych tam błędów.


\section{Literatura i materiały dodatkowe}

Materiałów jest mnóstwo. Oto parę z nich:
\begin{itemize}
    \item \emph{The Not So Short Introduction\ldots}, która posiada również tłumaczenie 
    w języku polskim.\\
    \url{http://www.ctan.org/tex-archive/info/lshort/english/lshort.pdf}

    \item Klasy stylu \texttt{memoir} posiadają bardzo wiele informacji o składzie tekstów
    anglosaskich oraz sposoby dostosowania \LaTeX{}a do własnych potrzeb.\\
    \url{http://www.ctan.org/tex-archive/macros/latex/contrib/memoir/memman.pdf}
    
    \item Nasza grupa dyskusyjna i repozytorium SVN są również dobrym miejscem aby zapytać
    (lub sprawdzić czy pytanie nie zostało już zadane).\\
    \url{https://ophelia.cs.put.poznan.pl/svn/put-latex/trunk}

    \item Dla łaknących więcej wiedzy o systemie \LaTeX{} podstawowym źródłem informacji
    jest książka Lamporta~\cite{Lamport:LDP85}. Prawdziwy \emph{hardcore} to oczywiście
    \emph{The \TeX{}book} profesora Knutha~\cite{Knuth:ct-a}.
\end{itemize}



% Bibliography (books, articles) starts here.
\scalebox{0.001}{
\cite{Lam98}
\cite{Lam01}}
\cite{CGR07}
\cite{AK08}
\cite{Urb03}
\bibliographystyle{alpha}{\raggedright\sloppy\small\bibliography{bibliography}}

% Colophon is a place where you should let others know about copyrights etc.
\ppcolophon

\end{document}
